\documentclass[10pt]{article}

\usepackage[margin=1.05in]{geometry}
\usepackage{amsmath,amssymb,mathtools}
\usepackage{booktabs}
\usepackage{enumitem}
\usepackage{xcolor}
\usepackage{microtype}
\usepackage[colorlinks=true,linkcolor=blue,citecolor=blue,urlcolor=blue]{hyperref}

\newcommand{\E}{\mathbb{E}}
\newcommand{\ind}[1]{\mathbf{1}!\left[#1\right]}
\newcommand{\sg}{\operatorname{sg}}
\newcommand{\Rho}{\rho}

\title{A Reach-First Successor Measure in the Forward--Backward Architecture}
\author{}
\date{}

\begin{document}

\maketitle

\section{Overview}

Forward--backward representation learning aims to learn the future occupancy of every
state under a family of policies:
\begin{equation}
M^{\pi_z}(s,a,X)
================

\E_{\pi_z}
\left[
\sum_{t\geq0}
\gamma^t
\ind{s_{t+1}\in X}
;\middle|;
s_0=s,;a_0=a
\right].
\label{eq:successor-measure}
\end{equation}

In continuous state spaces, FB represents the density of this measure relative to a
reference distribution $\rho$ as
\begin{equation}
m_z(s,a,x)
==========

\frac{dM^{\pi_z}(s,a,\cdot)}{d\rho}(x)
\approx
F(s,a,z)^\top B(x).
\label{eq:fb-factorization}
\end{equation}

The policy index $z$ and queried future state $x$ play different roles:
\begin{itemize}
\item $z$ selects the continuation policy $\pi_z$;
\item $x$ specifies the future state whose occupancy is queried;
\item $F(s,a,z)$ summarizes future behavior under $\pi_z$;
\item $B(x)$ provides coordinates for querying occupancy at $x$.
\end{itemize}

The proposed method preserves this architecture but changes the temporal-difference
objective applied to the bilinear prediction.

Stock FB learns a successor-measure density ratio using a least-squares density-estimation
objective.
The proposed curriculum instead makes the same inner product predict:
\begin{enumerate}
\item a bounded first-arrival value;
\item then a non-absorbing visitation value.
\end{enumerate}

The central change is therefore
\begin{equation}
\boxed{
\text{density-ratio estimation}
\quad\longrightarrow\quad
\text{ordinary Bellman regression on finite goal regions}.
}
\end{equation}

\section{Normal FB}

\subsection{Reward inference}

Suppose a downstream state reward is $r(x)$.
FB constructs a reward embedding
\begin{equation}
z_r
===

\E_{x\sim\rho}
\left[
r(x)B(x)
\right].
\label{eq:reward-embedding}
\end{equation}

Under the factorization in Equation~\eqref{eq:fb-factorization},
\begin{align}
F(s,a,z_r)^\top z_r
&=
F(s,a,z_r)^\top
\E_{x\sim\rho}
\left[
r(x)B(x)
\right]
\
&=
\E_{x\sim\rho}
\left[
r(x)
F(s,a,z_r)^\top B(x)
\right]
\
&\approx
Q_r^{\pi_{z_r}}(s,a).
\end{align}

The same latent vector $z_r$ is therefore used as:
\begin{itemize}
\item a representation of the reward;
\item a conditioning variable for the successor representation;
\item a conditioning variable for the actor.
\end{itemize}

The actor is trained to maximize
\begin{equation}
Q_z(s,a)
========

F(s,a,z)^\top z,
\label{eq:fb-q}
\end{equation}
so that
\begin{equation}
\pi_z(\cdot\mid s)
\approx
\arg\max_a F(s,a,z)^\top z.
\end{equation}

\subsection{Correspondence with the implementation}

The implementation contains the three FB components
\begin{align}
B(x)
&\longleftrightarrow
\texttt{_backward_map},
\
F(s,a,z)
&\longleftrightarrow
\texttt{_forward_map},
\
\pi(a\mid s,z)
&\longleftrightarrow
\texttt{_actor}.
\end{align}

For a batch
\begin{equation}
{(s_i,a_i,s_i')}_{i=1}^{n},
\end{equation}
the implementation samples one policy latent $z_i$ for each row and uses the batch of
successor states as query states:
\begin{equation}
x_j=s_j'.
\end{equation}

It computes
\begin{equation}
F_{hi}
======

F_h(s_i,a_i,z_i)
\in\mathbb{R}^d
\end{equation}
for each forward head $h$, and
\begin{equation}
B_j=B(x_j)\in\mathbb{R}^d.
\end{equation}

The matrix multiplication
\begin{equation}
M_{hij}
=======

F_h(s_i,a_i,z_i)^\top B(x_j)
\label{eq:batch-matrix}
\end{equation}
corresponds to
\begin{verbatim}
Fs = forward_map(left_enc, z, action)
B  = backward_map(goal)
Ms = torch.matmul(Fs, B.T)
\end{verbatim}

Each row $i$ fixes:
\begin{itemize}
\item a state-action pair $(s_i,a_i)$;
\item a continuation policy indexed by $z_i$.
\end{itemize}

Each column $j$ asks:
\begin{quote}
How much future occupancy does policy $\pi_{z_i}$ assign to query state $x_j$?
\end{quote}

Thus, the tensor already has the separated FB structure
\begin{equation}
(z_i,x_j)
\longmapsto
M^{\pi_{z_i}}(s_i,a_i,x_j).
\end{equation}

\section{The stock FB temporal-difference objective}

\subsection{Successor-density Bellman equation}

The successor-density ratio formally satisfies
\begin{equation}
m_z(s,a,x)
==========

\frac{\delta_{s'}(x)}{\rho(x)}
+
\gamma
\E_{a'\sim\pi_z(\cdot\mid s')}
\left[
m_z(s',a',x)
\right].
\label{eq:density-bellman}
\end{equation}

The immediate term
\begin{equation}
\frac{\delta_{s'}(x)}{\rho(x)}
\end{equation}
is a Dirac density ratio.
In a continuous state space, it is not an ordinary finite scalar target that can be
placed directly inside a squared Bellman residual.

FB therefore uses a least-squares density-ratio objective rather than directly regressing
Equation~\eqref{eq:density-bellman}.

\subsection{Bootstrapped prediction}

For each row, the actor supplies a continuation action
\begin{equation}
a_i'
\sim
\pi(\cdot\mid s_i',z_i).
\end{equation}

The target network predicts
\begin{equation}
\overline{M}_{ij}'
==================

\overline{F}(s_i',a_i',z_i)^\top
\overline{B}(x_j).
\end{equation}

Define the temporal-difference quantity
\begin{equation}
\Delta_{hij}
============

## M_{hij}

\gamma_i\overline{M}_{ij}'.
\label{eq:fb-diff}
\end{equation}

This corresponds to
\begin{verbatim}
diff = Ms - discount * target_M
\end{verbatim}

\subsection{Off-diagonal term}

For $i\neq j$, the query state $x_j$ is approximately independent of the transition in
row $i$.
These entries approximate samples from the reference distribution $\rho$.

The quadratic term is
\begin{equation}
\mathcal{L}_{\mathrm{off}}
==========================

\frac{1}{2}
\E_{i\neq j}
\left[
\Delta_{ij}^2
\right].
\label{eq:offdiag-loss}
\end{equation}

In the implementation,
\begin{verbatim}
fb_offdiag =
0.5 * (diff * off_diag).pow(2).sum() / off_diag_sum
\end{verbatim}

\subsection{Diagonal term}

On the diagonal,
\begin{equation}
x_i=s_i',
\end{equation}
so the query is the successor that was actually observed.

The linear matched-successor term is
\begin{equation}
\mathcal{L}_{\mathrm{diag}}
===========================

*

\E_i
\left[
\Delta_{ii}
\right].
\label{eq:diag-loss}
\end{equation}

In the implementation,
\begin{verbatim}
fb_diag =
-torch.diagonal(diff, dim1=1, dim2=2).mean()
\end{verbatim}

This term encourages the model to assign high occupancy to observed successor states.

The resulting FB loss is approximately
\begin{equation}
\mathcal{L}_{\mathrm{FB}}
=========================

\frac{1}{2}
\E_{x\sim\rho}
\left[
\left(
m_z(s,a,x)
----------

\gamma m_z(s',a',x)
\right)^2
\right]
-------

\E
\left[
m_z(s,a,s')
-----------

\gamma m_z(s',a',s')
\right]
+
\lambda_{\mathrm{ortho}}
\mathcal{L}_{\mathrm{ortho}}.
\label{eq:stock-fb-loss}
\end{equation}

The diagonal is therefore not trained against an explicit target such as $1$.
It enters through a negative linear term.

\section{The role of the backward-map regularizer}

The FB objective identifies the product
\begin{equation}
F(s,a,z)^\top B(x),
\end{equation}
but does not uniquely identify the individual scales and geometry of $F$ and $B$.

The implementation consequently adds an approximate whitening objective on the backward
embedding.

For a batch matrix
\begin{equation}
C_{ij}=B(x_i)^\top B(x_j),
\end{equation}
the regularizer is
\begin{equation}
\mathcal{L}_{\mathrm{ortho}}
============================

\frac{1}{2}
\E_{i\neq j}
\left[
C_{ij}^2
\right]
-------

\E_i[C_{ii}].
\label{eq:ortho-loss}
\end{equation}

The off-diagonal term discourages correlated embeddings, while the diagonal term prevents
the embeddings from collapsing to zero.

In the stock algorithm:
\begin{itemize}
\item the linear successor term pushes predicted measure values upward;
\item the quadratic term penalizes large density-ratio residuals;
\item the orthonormality objective constrains the geometry of $B$.
\end{itemize}

The manipulation experiments indicate that this balance is poorly conditioned:
orthonormality damages the useful reaching geometry, while removing it does not prevent
the density-based critic from inflating.

\section{The policy--measure feedback loop}

The actor uses
\begin{equation}
Q_z(s,a)
========

F(s,a,z)^\top z
\end{equation}
and is updated to increase this value.

At the same time, $F$ is trained to represent the successor measure generated by that
actor:
\begin{equation}
F(s,a,z)^\top B(x)
\approx
M^{\pi_z}(s,a,x).
\end{equation}

The complete dependency is
\begin{equation}
F
\longrightarrow
Q_z=F^\top z
\longrightarrow
\pi_z
\longrightarrow
M^{\pi_z}
\longrightarrow
F.
\end{equation}

The target measure changes as the actor changes, while the actor changes according to the
learned measure.
This circularity is one reason that the practical FB update is not simply ordinary
fixed-policy temporal-difference evaluation.

\section{Replacing density estimation with finite-region regression}

\subsection{Finite goal regions}

Let $c(s)\in\mathbb{R}^3$ denote the cube position in state $s$.
For a query state $x$, define the finite-region event
\begin{equation}
k_\epsilon(s,x)
===============

\ind{
|c(s)-c(x)|_2
\leq
\epsilon
}.
\label{eq:kernel}
\end{equation}

For the cube experiments,
\begin{equation}
\epsilon=4\text{ cm}.
\end{equation}

Instead of asking for the density of observing exactly $x$, the model asks whether the
next cube position lies inside a finite ball around $x$.

This replaces the singular immediate term
\begin{equation}
\frac{\delta_{s'}(x)}{\rho(x)}
\end{equation}
with an ordinary scalar:
\begin{equation}
k_\epsilon(s',x)\in{0,1}.
\end{equation}

Every element of the state-query matrix can now be trained with an explicit squared
Bellman residual.

\section{Phase 1: first-hitting values in the FB architecture}

\subsection{The represented object}

For a fixed policy index $z$ and query state $x$, define the first hitting time
\begin{equation}
T_x
===

\min
\left{
t\geq1:
k_\epsilon(s_t,x)=1
\right}.
\end{equation}

The discounted reaching value is
\begin{equation}
H^{\pi_z}(s,a,x)
================

\E_{\pi_z}
\left[
\gamma^{T_x-1}
\right].
\label{eq:hitting-value}
\end{equation}

We use the FB factorization
\begin{equation}
H_\theta(s,a,z,x)
=================

\frac{
F_\theta(s,a,z)^\top B_\omega(x)
}{
\sqrt{d}
}.
\label{eq:hitting-factorization}
\end{equation}

The $1/\sqrt{d}$ normalization matches the bilinear goal-conditioned critic used in the
successful baseline.

\subsection{Batch construction}

For a batch of transitions and query states,
\begin{equation}
{(s_i,a_i,s_i',z_i)}*{i=1}^{n},
\qquad
{x_j}*{j=1}^{q},
\end{equation}
define the hit matrix
\begin{equation}
K_{ij}
======

# k_\epsilon(s_i',x_j)

\ind{
|c(s_i')-c(x_j)|_2\leq\epsilon
}.
\label{eq:hit-matrix}
\end{equation}

The continuation action depends on the policy index but not on the query:
\begin{equation}
a_i'
\sim
\pi(\cdot\mid s_i',z_i).
\label{eq:policy-consistent-action}
\end{equation}

The same $a_i'$ is used to evaluate every query $x_j$ under policy $\pi_{z_i}$.

\subsection{Absorbing target}

The reaching target is
\begin{equation}
Y^{\mathrm{reach}}_{ij}
=======================

K_{ij}
+
\gamma_i
(1-K_{ij})
\overline{H}(s_i',a_i',z_i,x_j).
\label{eq:reaching-target}
\end{equation}

If the next state enters the query region, then $K_{ij}=1$ and the bootstrap terminates.

The loss is
\begin{equation}
\mathcal{L}_{\mathrm{reach}}
============================

\frac{1}{Hnq}
\sum_{h=1}^{H}
\sum_{i=1}^{n}
\sum_{j=1}^{q}
\left(
H_{h}(s_i,a_i,z_i,x_j)
----------------------

\sg\left[
Y^{\mathrm{reach}}_{ij}
\right]
\right)^2.
\label{eq:reaching-loss}
\end{equation}

There is:
\begin{itemize}
\item no separate diagonal loss;
\item no off-diagonal mask;
\item no linear matched-successor term.
\end{itemize}

Every matrix entry is an ordinary Bellman regression target.

When the queries are the batch successors, $x_j=s_j'$, the diagonal entries satisfy
\begin{equation}
K_{ii}=1.
\end{equation}

The diagonal is therefore an immediate positive example rather than a special linear
density-estimation term.

\subsection{Why the scale is anchored}

The fixed point satisfies
\begin{equation}
0
\leq
H^{\pi_z}(s,a,x)
\leq
1.
\end{equation}

An immediate hit has target value $1$.
The scalar event therefore anchors the scale of the bilinear product directly, removing
the main reason for using an auxiliary orthonormality constraint.

\section{Phase 2: non-absorbing visitation values}

\subsection{The represented object}

After learning the reaching representation, define the finite-region visitation value
\begin{equation}
V_\epsilon^{\pi_z}(s,a,x)
=========================

\E_{\pi_z}
\left[
\sum_{t\geq0}
\gamma^t
k_\epsilon(s_{t+1},x)
\right].
\label{eq:visitation-value}
\end{equation}

Use the same architecture:
\begin{equation}
V_\theta(s,a,z,x)
=================

\frac{
F_\theta(s,a,z)^\top B(x)
}{
\sqrt{d}
}.
\label{eq:visitation-factorization}
\end{equation}

At the beginning of phase 2:
\begin{enumerate}
\item initialize $F$, $B$, and the actor from phase 1;
\item freeze $B$;
\item warm-start $F$ from the reaching solution.
\end{enumerate}

\subsection{Non-absorbing target}

The phase-2 target is
\begin{equation}
Y^{\mathrm{visit}}_{ij}
=======================

K_{ij}
+
\gamma_i
\overline{V}(s_i',a_i',z_i,x_j).
\label{eq:visitation-target}
\end{equation}

The factor $(1-K_{ij})$ is removed.
Reaching the region produces one unit of occupancy, but the backup continues so that
dwelling and future revisits are also counted.

The loss is
\begin{equation}
\mathcal{L}_{\mathrm{visit}}
============================

\frac{1}{Hnq}
\sum_{h=1}^{H}
\sum_{i=1}^{n}
\sum_{j=1}^{q}
\left(
V_h(s_i,a_i,z_i,x_j)
--------------------

\sg\left[
Y^{\mathrm{visit}}_{ij}
\right]
\right)^2.
\label{eq:visitation-loss}
\end{equation}

The only mathematical change between the two phases is therefore
\begin{align}
Y^{\mathrm{reach}}
&=
K+\gamma(1-K)\overline{H}',
\
Y^{\mathrm{visit}}
&=
K+\gamma\overline{V}'.
\end{align}

Phase 1 terminates at first arrival.
Phase 2 counts every visit.

\section{Why the curriculum makes sense for play data}

Manipulation play trajectories contain substantial dwelling and recurrence.
After a cube reaches a region, the dataset may contain many nearly identical transitions
in which the cube remains there, followed by additional returns to the same region.

A direct successor-measure target is therefore numerically dominated by post-arrival
occupancy.

The policy-learning question, however, is largely determined before arrival:
\begin{quote}
Which action moves the cube toward the desired region quickly and reliably?
\end{quote}

The reaching value isolates this action-dependent signal:
\begin{equation}
H(s,a,z,x)
==========

\E[\gamma^{T_x-1}].
\end{equation}

The visitation value additionally contains dwelling and recurrence:
\begin{equation}
V(s,a,z,x)
==========

\E
\left[
\sum_{t\geq0}
\gamma^t
k_\epsilon(s_{t+1},x)
\right].
\end{equation}

For a small region, one may heuristically write
\begin{equation}
V(s,a,z,x)
\approx
H(s,a,z,x)R_z(x),
\label{eq:reach-recurrence}
\end{equation}
where $R_z(x)$ describes post-arrival recurrence.

The recurrent factor can be much larger than the reaching value.
Training directly on visitation therefore asks the model to discover a relatively small
action-dependent reaching signal underneath a large recurrence-dominated target.

The curriculum changes the order of learning:
\begin{enumerate}
\item first organize the representation according to reaching;
\item then add the repeated occupancy that successor measures are intended to count.
\end{enumerate}

The successor measure still performs long-horizon occupancy accounting, but recurrence
does not determine the representation before the control-relevant geometry has been
learned.

\section{Relation to the existing IQL critic mode}

The implementation already contains a related mode in which
\begin{equation}
Q(s,a,z)
========

F(s,a,z)^\top z
\end{equation}
is trained as a goal-reaching value using IQL.

With a task-independent forward map and bilinear value normalization, this becomes
approximately
\begin{equation}
Q(s,a,g)
========

\frac{
\phi(s,a)^\top\psi(g)
}{
\sqrt{d}
}.
\end{equation}

This existing mode implements a diagonal goal-reaching critic:
\begin{equation}
z_g=B(g),
\qquad
Q(s,a,g)
========

F(s,a,z_g)^\top B(g).
\end{equation}

It differs from the proposed matrix reaching objective in three ways:
\begin{enumerate}
\item it usually trains one selected goal per row rather than all state-query pairs;
\item it uses an expectile value function $V(s,z)$;
\item it does not necessarily learn $H(s,a,z,x)$ for independent $z$ and $x$.
\end{enumerate}

Most importantly, switching from the IQL critic back to the current measure mode is not
phase 2.
The current measure mode restores the original density-ratio loss:
\begin{equation}
\mathcal{L}*{\mathrm{off}}
+
\mathcal{L}*{\mathrm{diag}}
+
\lambda\mathcal{L}_{\mathrm{ortho}}.
\end{equation}

The proposed phase 2 instead requires a new squared-residual visitation mode.

A useful implementation distinction is therefore
\begin{verbatim}
critic_mode = "measure"       # stock FB density-ratio loss
critic_mode = "reach_kernel"  # absorbing finite-region TD
critic_mode = "visit_kernel"  # non-absorbing finite-region TD
\end{verbatim}

\section{Diagonal and fully separated versions}

\subsection{Diagonal goal-conditioned curriculum}

The simplest version ties the policy and query to the same goal:
\begin{equation}
z_g=B(g),
\end{equation}
and learns
\begin{equation}
H^{\pi_g}(s,a,g)
\end{equation}
followed by
\begin{equation}
V^{\pi_g}(s,a,g).
\end{equation}

Equivalently, it learns
\begin{equation}
g
\longmapsto
V^{\pi_g}
\left(
s,a,B_\epsilon(g)
\right).
\end{equation}

This supports:
\begin{itemize}
\item goal-conditioned control;
\item nearest-goal retrieval;
\item generalized policy improvement over a small candidate set.
\end{itemize}

It does not provide the full successor measure needed to evaluate arbitrary rewards
because the policy and occupancy query are tied to the same goal.

\subsection{Fully separated FB curriculum}

The full version preserves independent policy and query variables:
\begin{equation}
H(s,a,z,x)
==========

F(s,a,z)^\top B(x)
\end{equation}
in phase 1, followed by
\begin{equation}
V(s,a,z,x)
==========

F(s,a,z)^\top B(x)
\end{equation}
in phase 2.

The batch tensor in Equation~\eqref{eq:batch-matrix} already has this structure.

This version is required for reward inference:
\begin{equation}
z_r
===

\E_{x\sim\rho}
\left[
r(x)B(x)
\right],
\end{equation}
followed by direct policy execution with
\begin{equation}
\pi(a\mid s,z_r).
\end{equation}

\section{Choosing policy latents during phase 1}

For a goal-derived latent
\begin{equation}
z_g=B(g),
\end{equation}
the policy semantics are clear:
$\pi_{z_g}$ should reach $g$.

For an arbitrary random direction $z$, there is no unique first-arrival event naturally
associated with that vector.
A linear combination of first-hitting values is not generally the value of the
corresponding linear combination of rewards.

A conservative curriculum is therefore:
\begin{enumerate}
\item train phase 1 primarily with goal-derived latents;
\item introduce reward-mixture or random latents once the non-absorbing visitation
representation is available.
\end{enumerate}

One possible schedule is
\begin{align}
p_{\mathrm{goal}}
&:
1.0\longrightarrow0.5,
\
p_{\mathrm{random}}
&:
0.0\longrightarrow0.5.
\end{align}

Random-$z$ reaching training should otherwise be described as a surrogate objective rather
than exact reward-conditioned policy evaluation.

\section{Recommended phase-2 ablations}

The cleanest phase-2 experiment freezes both:
\begin{itemize}
\item the backward representation $B$;
\item the phase-1 actor $\pi_z$.
\end{itemize}

Then phase 2 becomes fixed-policy visitation evaluation:
\begin{equation}
V^{\pi_z}
=========

K_\epsilon
+
\gamma P^{\pi_z}V^{\pi_z}.
\end{equation}

For a fixed policy, this is an ordinary discounted Bellman operator.

Three useful configurations are:
\begin{center}
\begin{tabular}{lll}
\toprule
\textbf{Backward map} & \textbf{Actor} & \textbf{Purpose} \
\midrule
Frozen & Frozen & Isolate visitation-value learning \
Frozen & Updated & Test visitation-based policy improvement \
Updated & Updated & Full end-to-end curriculum \
\bottomrule
\end{tabular}
\end{center}

The first configuration most directly tests whether a reaching representation can be
converted into an accurate visitation representation.

Updating the actor during phase 2 reintroduces a policy--measure feedback loop:
the actor changes according to the visitation value while the visitation target changes
according to the actor.

\section{Implementation inside \texttt{update_fb}}

The existing forward computation can be retained:
\begin{verbatim}
Fs = forward_map(left_enc, z, action)
B  = backward_map(query_obs)
values = torch.matmul(Fs, B.T) * value_scale
\end{verbatim}

The target forward computation can also be retained:
\begin{verbatim}
next_action = actor(next_obs, z)
target_Fs   = target_forward(next_obs, z, next_action)
target_B    = target_backward(query_obs)
target_values = torch.matmul(target_Fs, target_B.T) * value_scale
\end{verbatim}

The new component is the hit matrix:
\begin{verbatim}
hit = (
torch.cdist(next_cube.float(), query_cube.float())
<= epsilon
).float()
\end{verbatim}

For phase 1:
\begin{verbatim}
continuation = 1.0 - hit
target = hit + discount * continuation * target_values
\end{verbatim}

For phase 2:
\begin{verbatim}
continuation = torch.ones_like(hit)
target = hit + discount * continuation * target_values
\end{verbatim}

The representation loss becomes
\begin{verbatim}
loss = (
values - target.detach().unsqueeze(0)
).pow(2).mean()
\end{verbatim}

This replaces:
\begin{verbatim}
diff = Ms - discount * target_M
fb_offdiag, fb_diag = fb_successor_terms(...)
fb_loss = fb_offdiag + fb_diag + ortho_loss
\end{verbatim}

The proposed kernel modes therefore do not require:
\begin{itemize}
\item the off-diagonal mask;
\item the diagonal linear term;
\item density-ratio orthonormalization.
\end{itemize}

\section{Reward inference and the kernel interpretation}

For the fully separated visitation representation, the existing reward-inference API can
still compute
\begin{equation}
z_r
===

\operatorname{Proj}
\left(
\sum_j r_jB(x_j)
\right).
\end{equation}

However, the learned quantity is now occupancy of finite balls:
\begin{equation}
V_\epsilon^{\pi_z}(s,a,x)
=========================

\E_{\pi_z}
\left[
\sum_{t\geq0}
\gamma^t
k_\epsilon(s_{t+1},x)
\right],
\end{equation}
rather than the exact density ratio
\begin{equation}
\frac{dM^{\pi_z}}{d\rho}(x).
\end{equation}

Consequently, reward inference naturally corresponds to a kernel-smoothed reward.
Overlapping balls and a non-uniform replay distribution may require:
\begin{itemize}
\item kernel normalization;
\item density correction;
\item or an explicit interpretation of downstream rewards in the same kernel basis.
\end{itemize}

This issue does not affect the reach-first curriculum itself, but it matters for the claim
of exact arbitrary-reward integration.

\section{Algorithm summary}

\paragraph{Stock FB.}
Represent
\begin{equation}
m_z(s,a,x)
==========

F(s,a,z)^\top B(x)
\end{equation}
and train it using
\begin{equation}
\frac{1}{2}
\E_\rho
\left[
(m-\gamma m')^2
\right]
-------

\E_{\mathrm{successor}}
[m-\gamma m']
+
\lambda\mathcal{L}_{\mathrm{ortho}}.
\end{equation}

\paragraph{Phase 1: reach-first FB.}
Represent
\begin{equation}
H_z(s,a,x)
==========

\frac{F(s,a,z)^\top B(x)}{\sqrt d}
\end{equation}
and train it using
\begin{equation}
H
\leftarrow
k_\epsilon(s',x)
+
\gamma
\bigl(1-k_\epsilon(s',x)\bigr)
H'.
\end{equation}

\paragraph{Phase 2: visitation FB.}
Represent
\begin{equation}
V_z(s,a,x)
==========

\frac{F(s,a,z)^\top B(x)}{\sqrt d}
\end{equation}
and train it using
\begin{equation}
V
\leftarrow
k_\epsilon(s',x)
+
\gamma V'.
\end{equation}

The network structure remains unchanged.
The principal modification is the loss applied to the matrix
\begin{equation}
F(s_i,a_i,z_i)^\top B(x_j).
\end{equation}

Stock FB interprets this matrix through density-ratio estimation.
The proposed curriculum first interprets it as first-hitting values and then as
finite-region visitation values.

\section{Central interpretation}

The method divides responsibilities between temporal-difference learning and the
successor representation.

Temporal-difference learning first solves the control problem:
\begin{quote}
Which actions reach a queried region quickly and reliably?
\end{quote}

The successor representation then solves the occupancy problem:
\begin{quote}
Once the reaching geometry is established, how much discounted time does each policy
spend in each queried region?
\end{quote}

This division is well matched to manipulation play data.
Repeated observations, dwelling, and revisitation dominate raw occupancy magnitude, but
they need not dominate the learning of action-relevant geometry.

The reach-first phase isolates the bounded action signal.
The visitation phase subsequently restores the long-horizon occupancy semantics required
by successor measures.

\end{document}
